% Research continuation, pinned against ProveIt 4c173c06cc32c9cea554b39be1837ad2ae897fc1.
% Standalone SECTION fragment. Requires amsmath,amssymb,amsthm,booktabs,hyperref.
% The root article supplies theorem,lemma,proposition,corollary,remark environments.

\section{Polynomial separation and a larger uniform Lerch zero regime}
\label{polylerch:sec:main}

The incoming effective-saturation report~\cite{incoming-uniform} obtains an explicit separation
bound $(3n)^{-(n-2)}$ for the reciprocal-gamma Appell polynomials and
a sufficient derivative order of size $n^{2n+O(n)}$.  Two refinements
give a polynomial sufficient order. First, a factorial diagonal
multiplier converts a rational initial segment of the gamma product
into a polynomial with quantitatively separated negative roots.
Second, summing reciprocal distances from all roots replaces the
coarser product estimate that charges every root the nearest distance.
The resulting sufficient order is $O(n^6\log^2 n)$, uniformly in the
entire Lerch deformation interval. This is a new consequence of
classical multiplier and mesh theory; the underlying preservation
theorems are not new claims.

Write
\[
 \frac{e^{xz}}{\Gamma(1+z)}
   =\sum_{n\ge0}Q_n(x)\frac{z^n}{n!},\qquad
 W_\rho(t)=\frac1{1-\rho e^{-t}}.
\]
For $n\ge0$, $k\ge1$, and $0\le\rho\le1$, retain the manuscript's
signed family
\begin{equation}\label{polylerch:eq:F}
 F_{n,k}^{\rho}(a)=(-1)^{n+k}D_{n,k}^{\rho}(a)
 =\int_0^\infty t^k e^{-at}W_\rho(t)Q_n(\log t)\,dt.
\end{equation}
Thus $D_{n,k}^{1}=\gamma_n^{(k)}$, and for $\rho<1$ it is the
$k$th $a$ derivative of
$\sum_{m\ge0}\rho^m\log^n(a+m)/(a+m)$.
Here the Stieltjes convention is
$\zeta(s,a)=(s-1)^{-1}+\sum_{n\ge0}(-1)^n\gamma_n(a)(s-1)^n/n!$.
The established Appell--Laplace formula is
\cite[\texttt{zeros:thm:laplace}]{baseline}; it also follows by differentiating
$t^z/\Gamma(1+z)$ under the defining gamma integral. We recall the
variation proof below so that the completeness of the zero list does
not depend on a numerical root search.

\subsection{A factorial multiplier and a rational separation bound}

For a polynomial whose nonzero roots all have one sign, order their
absolute values as $0<r_1\le\cdots\le r_d$ and define its logarithmic
mesh to be $\min_{i<d}r_{i+1}/r_i$. Ordinary mesh means the minimum
additive distance between consecutive roots, counting multiplicity.

We use the following two classical preservation statements. A
real-rootedness-preserving operator commuting with translations does
not reduce ordinary mesh. A real-rootedness-preserving diagonal
operator does not reduce logarithmic mesh on polynomials with all
roots of one sign. The second statement is the theorem of Katkova,
Shapiro and Vishnyakova~\cite{polylerch:KSV}; the first includes the
Riesz theorem for $1+c\partial_x$. Both follow from the interlacing
pencil criterion: apply the operator to every real linear combination
of $p(x),p(x+h)$, or $p(x),p(qx)$, respectively. Interlacing of the
image pair gives the asserted mesh inequality. Here $h$ is below
the original additive mesh and $1<q$ below the original logarithmic
mesh; passage to the supremum includes equality.

\begin{lemma}[A concrete diagonal multiplier]\label{polylerch:lem:multiplier}
The operator
\[
 \mathcal M\left(\sum_{r=0}^d a_rx^r\right)
   =\sum_{r=0}^d\frac{a_r}{(r+1)!}x^r
\]
preserves real-rootedness. If $p$ has strictly negative roots and positive
leading coefficient, then $\mathcal Mp$ has strictly negative roots
and logarithmic mesh at least that of $p$.
\end{lemma}
\begin{proof}
Put $E=x\partial_x$ and, for integers $N\ge1$, define
\[
 \mathcal M_Np(x)
 =\left[\prod_{j=2}^{N+1}(1+E/j)\right]p(x/N).
\]
Each factor preserves real-rootedness. Indeed, for positive integral $j$,
\[
 (1+E/j)p(x)=\frac{x^{1-j}}j\frac{d}{dx}\bigl(x^jp(x)\bigr).
\]
Rolle's theorem makes the derivative real-rooted; the division merely
removes $j-1$ roots at zero. Positive dilation also preserves
real-rootedness. On the monomial $x^r$ the multiplier is
\[
 N^{-r}\prod_{j=2}^{N+1}(1+r/j)
 =\frac1{(r+1)!}\prod_{h=2}^{r+1}(1+h/N)
 \longrightarrow\frac1{(r+1)!}.
\]
The empty product at $r=0$ equals one. Thus $\mathcal M_Np$ converges
coefficientwise to $\mathcal Mp$. Its degree does not drop, and closure
of real-rootedness proves the first assertion. The coefficients of a
polynomial with negative roots and positive leading coefficient are
strictly positive. Their images remain strictly positive, so the
real roots of $\mathcal Mp$ cannot be nonnegative. The classical
diagonal mesh statement now applies. In particular, a logarithmic
mesh strictly greater than one proves simplicity of the image.
\end{proof}

\begin{theorem}[Polynomial Appell separation]\label{polylerch:thm:mesh}
For every $n\ge3$, the polynomial $Q_n$ has $n$ distinct real roots and
\begin{equation}\label{polylerch:eq:delta}
 \operatorname{mesh}(Q_n)\ge\Delta_n,
 \qquad \Delta_n=\frac4{n(n-1)(n-2)}.
\end{equation}
For $n=2$, one may take $\Delta_2=2$.
\end{theorem}
\begin{proof}
Set
\[
 B_n(x)=\prod_{j=1}^{n-1}(1+\partial_x/j)x^n,
 \qquad P_n(x)=\prod_{j=1}^{n-1}(1+jx).
\]
Complementation of elementary symmetric functions gives the exact
identity
\begin{equation}\label{polylerch:eq:Bidentity}
 B_n(x)=n x R_n(x),\qquad R_n(x)=\mathcal MP_n(x).
\end{equation}
For explicit verification, if $e_r$ denotes the degree-$r$ elementary
symmetric function of $1,2,\ldots,n-1$, then
\[
 R_n(x)=\sum_{r=0}^{n-1}\frac{e_r}{(r+1)!}x^r,
 \qquad
 e_{n-1-r}(1,1/2,\ldots,1/(n-1))=\frac{e_r}{(n-1)!}.
\]

The negative roots of $P_n$ have absolute values
$1/(n-1),1/(n-2),\ldots,1$. Its logarithmic mesh is
$(n-1)/(n-2)$. By Lemma~\ref{polylerch:lem:multiplier}, the roots of
$R_n$ can be written as $-r_1,\ldots,-r_{n-1}$ with
\[
 0<r_1<\cdots<r_{n-1},\qquad
 \frac{r_{i+1}}{r_i}\ge\frac{n-1}{n-2}.
\]
The coefficient ratio supplies an absolute scale:
\[
 \sum_{i=1}^{n-1}\frac1{r_i}
 =\frac{R_n'(0)}{R_n(0)}
 =\frac12\sum_{j=1}^{n-1}j=\frac{n(n-1)}4.
\]
Consequently $r_1\ge4/[n(n-1)]$. Every consecutive nonzero-root
gap is at least $r_1/(n-2)$, and the remaining gap, from $-r_1$ to
the single zero at the origin in \eqref{polylerch:eq:Bidentity}, is
$r_1$ itself. This proves $\operatorname{mesh}(B_n)\ge\Delta_n$.
For $n=2$, $B_2=x(x+2)$ directly gives mesh two.

For $N\ge n-1$ define the finite reciprocal-gamma Appell polynomial
\[
 Q_{n,N}(x)=\left[\prod_{j=1}^{N}(1+\partial_x/j)\right]
                  (x+\gamma-H_N)^n.
\]
The first $n-1$ factors give a translation of $B_n$. Each remaining
factor preserves real-rootedness and, by Riesz's theorem, does not
decrease additive mesh. Hence $\operatorname{mesh}(Q_{n,N})\ge\Delta_n$.
The Weierstrass product
\[
 \frac1{\Gamma(1+z)}=e^{\gamma z}
             \prod_{j\ge1}(1+z/j)e^{-z/j}
\]
converges locally uniformly, so $Q_{n,N}$ converges coefficientwise
to $Q_n$. Continuity of the ordered roots of monic real-rooted
polynomials preserves the positive gap inequalities in the limit.
This proves both the stated bound and simplicity.
\end{proof}

The zero at the origin belongs only to the auxiliary $B_n$. No
logarithmic mesh is assigned to that zero: the multiplicative argument
is applied to $R_n$, and the final additive gap is handled separately.
This distinction is essential to the proof.

\subsection{The inherited global zero bound, with multiplicity}

\begin{lemma}[Laplace variation]\label{polylerch:lem:variation}
For $n\ge0$, $k\ge1$, and $0\le\rho\le1$, the function
$F_{n,k}^{\rho}$ has at most $n$ positive zeros counted with multiplicity.
\end{lemma}
\begin{proof}
Here is the elementary argument from the manuscript
\cite[\texttt{zeros:lem:variation}]{baseline}. If a nonzero real
kernel $g$ has $m$ sign changes and its exponentially weighted
polynomial moments exist, its Laplace transform $G$ has at most
$m$ positive zeros. For $m=0$ the transform has a strict constant
sign. For $m>0$, choose a sign-change point $t_0$.
The kernel $(t_0-t)g(t)$ has $m-1$ sign changes, and its transform is
\[
 (\partial_a+t_0)G(a)
   =e^{-t_0a}\frac{d}{da}\bigl(e^{t_0a}G(a)\bigr).
\]
Induction bounds its zeros by $m-1$ and precludes an identically
zero transform. Any finite collection of zeros of $G$ of total
multiplicity $M$ forces at least $M-1$ zeros of the displayed
transformed derivative, by Rolle's theorem with multiplicities.
Thus $M\le m$, completing induction, including nonvanishing.

The factor $t^kW_\rho(t)$ is strictly positive. By
Theorem~\ref{polylerch:thm:mesh}, $Q_n(\log t)$ has exactly $n$
simple sign changes; the cases $n=0,1$ are immediate.
Moreover $W_\rho(t)\le W_1(t)\le1+t^{-1}$, so $k\ge1$ gives
integrability at zero and the exponential controls infinity.
These bounds justify all differentiations used in the induction.
\end{proof}

\subsection{Sign control using the whole root configuration}

Let $Y_k$ have gamma shape $k+1$ and rate $k$. For a real number $s$
that is not a root of $Q_n$, set
\[
 S_n(s)=\sum_{j=1}^n\frac1{|s-x_{n,j}|},
 \qquad m(s)=\left\lceil1+S_n(s)\right\rceil.
\]
Here $x_{n,1}>\cdots>x_{n,n}$ are the roots of $Q_n$.

\begin{lemma}[Configuration-sensitive sign transfer]
\label{polylerch:lem:transfer}
If $k\ge16m(s)^2$, then uniformly for $0\le\rho\le1$,
\begin{equation}\label{polylerch:eq:transfer}
 \left|
 \frac{(ke^{-s})^{k+1}F_{n,k}^{\rho}(ke^{-s})}
 {k!W_\rho(e^s)Q_n(s)}-1\right|
 \le\sqrt{2/3}<1.
\end{equation}
Thus $F_{n,k}^{\rho}(ke^{-s})$ and $Q_n(s)$ have the same strict sign.
\end{lemma}
\begin{proof}
The exact gamma average obtained by setting $t=e^s y$ in
\eqref{polylerch:eq:F} is
\[
 \frac{(ke^{-s})^{k+1}F_{n,k}^{\rho}(ke^{-s})}
 {k!W_\rho(e^s)Q_n(s)}
 =\mathbb E\left[
 \frac{W_\rho(e^sY_k)}{W_\rho(e^s)}
 \frac{Q_n(s+\log Y_k)}{Q_n(s)}\right].
\]
For real $u$, factorization over the roots yields
\[
 \left|\frac{Q_n(s+u)}{Q_n(s)}-1\right|
 \le\prod_j\left(1+\frac{|u|}{|s-x_{n,j}|}\right)-1
 \le e^{S_n(s)|u|}-1.
\]
Also
\[
 \left|\frac{d\log W_\rho(t)}{d\log t}\right|
 =\frac{t\rho}{e^t-\rho}\le1,
\]
since $e^t-\rho\ge\rho t$. Therefore, writing
$R=W_\rho(e^{s+u})/W_\rho(e^s)$, one has
$R\le e^{|u|}$ and $|R-1|\le e^{|u|}-1$.
The absolute difference of the full integrand from one is at most
$e^{(1+S_n(s))|u|}-1\le e^{m(s)|u|}-1$.

For completeness, the required gamma moment estimate is
\[
 \mathbb E(e^{m|\log Y_k|}-1)
 \le\{2(e^{4m^2/k}-1)\}^{1/2}\le\sqrt{2/3}
 \quad(k\ge16m^2).
\]
Indeed, the squared integrand is bounded by
$(Y_k^m-1)^2+(Y_k^{-m}-1)^2$ and Cauchy--Schwarz applies. The exact
integer moments are
\[
 \mathbb EY_k^r=\prod_{j=1}^{r}(1+j/k),\qquad
 \mathbb EY_k^{-r}=\prod_{j=0}^{r-1}(1-j/k)^{-1}.
\]
Both moments with $r=m$ are at least one. For $r=2m$, their
logarithms are at most $3m^2/k$ and $4m^2/k$, respectively, using
$\log(1+v)\le v$ and $-\log(1-v)\le2v$ for $0\le v\le1/2$.
The inverse moments exist because $2m\le k$. This gives the first
bound. Finally $4m^2/k\le1/4$ and
$e^{1/4}\le(1-1/4)^{-1}=4/3$ give the second. Taking expectations
proves \eqref{polylerch:eq:transfer} and also verifies absolute
integrability of every expression used.
\end{proof}

\subsection{The polynomial threshold and the simultaneous zero limit}

\begin{theorem}[Uniform saturation at polynomial derivative order]
\label{polylerch:thm:saturation}
For $n\ge3$ define the explicit integer
\begin{equation}\label{polylerch:eq:threshold}
 \mathcal K_n=
 \left\lceil\bigl(16+6n(n-1)(n-2)H_{n-1}\bigr)^2\right\rceil.
\end{equation}
For every integer $k\ge\mathcal K_n$ and every $0\le\rho\le1$,
the function $D_{n,k}^{\rho}$ has exactly $n$ positive zeros, all simple.
Numbering them increasingly as $a_{n,1,k}^{\rho},\ldots,a_{n,n,k}^{\rho}$
and putting $c_{n,j}=e^{-x_{n,j}}$, one has
\begin{equation}\label{polylerch:eq:location}
 \left|\log\frac{a_{n,j,k}^{\rho}}{kc_{n,j}}\right|
 <\frac8{\sqrt k}\qquad(1\le j\le n).
\end{equation}
In particular,
\[
 \mathcal K_n=36n^6\log^2 n\,(1+o(1)),\qquad
 \mathcal K_n\le36n^4(n-1)^2(n-2)^2\le36n^8.
\]
The exact two-zero theorem remains stronger at $n=2$; for the stated
location estimate alone, the same proof gives the sufficient value
$k\ge784$ with $\Delta_2=2$.
\end{theorem}
\begin{proof}
Let $\delta=\Delta_n$ and $\eta=8/\sqrt k$. The threshold condition is
precisely
\[
 \sqrt k\ge16+\frac{24H_{n-1}}\delta.
\]
It implies $\eta<\delta/3$, because $H_{n-1}\ge1$. The intervals
$[x_{n,j}-\eta,x_{n,j}+\eta]$ are disjoint. At either endpoint
$s=x_{n,j}\pm\eta$, the reciprocal-distance sum obeys
\begin{align*}
 S_n(s)&\le\frac1\eta+
       \frac3{2\delta}(H_{j-1}+H_{n-j})\\
 &\le\frac1\eta+\frac{3H_{n-1}}\delta.
\end{align*}
For a root $h$ positions away, this uses
$|s-x_{n,j\pm h}|\ge h\delta-\eta\ge2h\delta/3$.
Consequently
\[
 m(s)\le2+S_n(s)
 \le2+\frac{\sqrt k}8+\frac{3H_{n-1}}\delta
 \le\frac{\sqrt k}4.
\]
Lemma~\ref{polylerch:lem:transfer} gives opposite endpoint signs for
$F_{n,k}^{\rho}(ke^{-s})$ on each interval, uniformly in $\rho$.
There is therefore a positive zero in each of the $n$ corresponding
disjoint $a$ intervals. Lemma~\ref{polylerch:lem:variation}, which counts
multiplicities, makes this list complete and every zero simple. The
interval endpoints give \eqref{polylerch:eq:location}.

The asymptotic size follows from $H_{n-1}=\log n+O(1)$. For the
displayed elementary upper bound, use $H_{n-1}\le n-1$ and
$16\le6n(n-1)(n-2)$ for $n\ge3$. They give
$16+6n(n-1)(n-2)H_{n-1}\le6n^2(n-1)(n-2)$;
the square on the right is already integral, so the ceiling preserves
the inequality. For $n=2$, the threshold becomes $(16+24/2)^2=784$.
\end{proof}

\begin{remark}[Sharper finite thresholds from certified auxiliary roots]
The proof uses only a lower bound $\delta$ on the Appell mesh. Thus every
proved $\delta>0$ gives the sufficient order
$\lceil(16+24H_{n-1}/\delta)^2\rceil$, with the same logarithmic
location error $8/\sqrt k$. Exact rational isolation of all roots of
$B_n$ supplies such a $\delta$ by Riesz's theorem. The accompanying
script verifies the polynomial identity, isolates all 135 auxiliary
roots for $2\le n\le16$, and validates the additive and logarithmic
mesh inequalities with rational brackets. The following finite
thresholds use those certified gaps:
\[
\begin{array}{c|rrrrrrrr}
n&3&4&5&6&8&10&12&16\\\hline
\text{sufficient }k&3629&12996&36442&85413&327915&926207&2152710&8065605.
\end{array}
\]
They remain sufficient bounds, not least thresholds. No finite root
calculation is needed for the all-index theorem or its growing-index
corollary. The exact audit also checks 90 finite-multiplier coefficient
identities and seven gamma-moment bounds.
\end{remark}

\begin{corollary}[A power-sized growing-index regime]
\label{polylerch:cor:growing}
Fix $0<c<1$. For all sufficiently large integers $k$, simultaneously for
\[
 3\le n\le c\frac{k^{1/6}}{(\log k)^{1/3}},\qquad
 1\le j\le n,\qquad0\le\rho\le1,
\]
the function $D_{n,k}^{\rho}$ has exactly $n$ simple positive zeros and
\[
 \frac{a_{n,j,k}^{\rho}}{k e^{-x_{n,j}}}=1+O(k^{-1/2}).
\]
The implied constant can be chosen independently of $n,j,\rho$.
\end{corollary}
\begin{proof}
The defining quantity in \eqref{polylerch:eq:threshold} is increasing
in $n$. At the largest permitted index its ratio to $k$ has limsup
at most $c^6<1$: substitute
$n\sim c k^{1/6}(\log k)^{-1/3}$ into
$36n^6\log^2 n$, noting $\log n\sim\tfrac16\log k$.
Hence $k\ge\mathcal K_n$ throughout the range for all sufficiently
large $k$. Theorem~\ref{polylerch:thm:saturation} applies.
Exponentiation of the uniform logarithmic bound gives the conclusion;
for example, for $k\ge64$ its relative error is less than
$8e/\sqrt k$.
\end{proof}

\subsection{What has and has not been resolved}

This argument strengthens the effective growing-index theorem from a
logarithmic index range to a power-sized one. It gives a rigorous
polynomial answer to the incoming report's question about substantially
smaller explicit saturation orders. It does not determine the least
uniform threshold, and the bound on the Appell mesh is not claimed
optimal. In particular, the elementary endpoint still forces every
uniform simplicity threshold to be at least $n-1$.
Indeed, at $\rho=0$, writing $x=\log a$ gives
\[
 \frac{d^k}{da^k}\frac{\log^n a}{a}
   =a^{-k-1}\prod_{j=1}^k(\partial_x-j)x^n.
\]
For $k\le n-2$, the right-hand polynomial is divisible by
$x^{n-k}$, so its positive zero at $a=1$ is multiple.

Further precise targets are: improve the $n^{-3}$ Appell mesh bound;
determine whether the least uniform simplicity threshold is $n-1$;
replace the absolute log-moment estimate by a signed concentration
estimate to improve the $k^{-1/2}$ location error; and extend the
configuration-sensitive argument to all-order offsets with a growing
index. The lower-branch shape problem at index two belongs to a separate
line of results already developed in the source reports. We do not
count its previously established cases as new results here.

% Bibliography item for the parent article:
% \bibitem{polylerch:KSV}
% O. Katkova, B. Shapiro and A. Vishnyakova,
% Multiplier sequences and logarithmic mesh,
% C. R. Math. Acad. Sci. Paris 349 (2011), nos. 1--2, 35--38.
% doi:10.1016/j.crma.2010.11.031; arXiv:1010.6052.
% Primary publisher page:
% https://comptes-rendus.academie-sciences.fr/mathematique/articles/10.1016/j.crma.2010.11.031/
% arXiv primary full text: https://arxiv.org/pdf/1010.6052
